\documentclass[12pt]{article}
\usepackage{amsmath,amssymb,amsthm,bm,mathtools}
\usepackage[margin=1in]{geometry}
\usepackage[skip=0.6em]{parskip}
\usepackage[hidelinks]{hyperref}

\newtheorem{theorem}{Theorem}
\newtheorem{proposition}[theorem]{Proposition}
\newtheorem{lemma}[theorem]{Lemma}
\newtheorem{corollary}[theorem]{Corollary}
\theoremstyle{definition}
\newtheorem{definition}[theorem]{Definition}
\theoremstyle{remark}
\newtheorem{remark}[theorem]{Remark}
\newtheorem{conjecture}[theorem]{Conjecture}

\newcommand{\E}{\mathrm E}
\newcommand{\Pp}{\mathrm P}
\newcommand{\R}{\mathbf R}
\newcommand{\cC}{\mathcal C}
\newcommand{\cR}{\mathcal R}
\newcommand{\cF}{\mathcal F}
\newcommand{\cB}{\mathcal B}
\newcommand{\1}{\mathbf 1}

\begin{document}

\begin{center}
{\Large The general core--redundant Gaussian experiment}\\[4pt]
{\large Feasible support, first-moment rigidity, and gated domination}\\[4pt]
{\small Ivan A. Canay :: August 28, 2026}
\end{center}

\medskip

This note studies admissibility in a Gaussian experiment with an arbitrary
number of core and redundant moments and an unrestricted positive-definite
correlation matrix.  The main distinction is between models with one core
moment and models with at least two.  With one core moment, dominance can
produce a nonzero first-moment transfer and calibrated gates can improve some
tests.  With at least two core moments, every weak dominator of an exact-size
test must match its full first-moment vector at the origin.  This identity
implies that feasible support plus one arbitrary supporting halfspace is
sufficient for admissibility.

The results below are self-contained.  They do not provide a full
characterization when there are at least two core moments.  The first open
case is a region requiring two essential infeasible supporting halfspaces.

\section{The experiment and the decision problem}

Let
\[
X\sim N_d(\mu,\Omega),
\]
where $\Omega$ is a positive-definite correlation matrix.  Partition the
indices into two nonempty sets: the core indices $\cC$ and the redundant
indices $\cR$.  The maintained mean space is
\begin{equation}\label{eq:Lambda}
\Lambda
:=
\left\{\mu\in\R^d:
\mu_r\le \max_{c\in\cC}\mu_c
\text{ for every }r\in\cR
\right\}.
\end{equation}
The null and alternative are
\begin{align}
\Lambda_0
&:=\Lambda\cap(-\infty,0]^d,\label{eq:null}\\
\Lambda_1
&:=\Lambda\setminus\Lambda_0
=\left\{\mu\in\Lambda:
\max_{c\in\cC}\mu_c>0\right\}.\label{eq:alternative}
\end{align}
Thus a redundant coordinate can be positive only if some core coordinate is
positive.  The restriction is imposed on standardized local means; it is not
implied by the corresponding ordering of unstandardized drifts unless their
standard deviations agree.

For a Borel acceptance region $A$, write
\[
\phi_A(x):=\1\{x\notin A\}.
\]
A possibly randomized test is a Borel function $\psi:\R^d\to[0,1]$.  It is
level $\alpha$ if
\[
\sup_{\mu\in\Lambda_0}\E_\mu\psi(X)\le\alpha.
\]
It weakly dominates $\phi_A$ if
\[
\E_\mu\psi(X)\ge\E_\mu\phi_A(X)
\qquad\text{for every }\mu\in\Lambda_1,
\]
and it strictly dominates $\phi_A$ if the inequality is strict for at least
one alternative.  The test $\phi_A$ is $\alpha$-admissible if it is level
$\alpha$ and has no strict dominator.

We call $A$ \emph{nice} if it is closed, convex, coordinatewise lower, and
\begin{equation}\label{eq:exact-size}
\Pp_0(X\in A)=1-\alpha.
\end{equation}
Coordinatewise lowerness and $\Lambda_0\subseteq(-\infty,0]^d$ imply that
$\phi_A$ is level $\alpha$, with size attained at the origin.

\section{The dual-cone dichotomy}

For a cone $K\subseteq\R^d$, define
\[
K^*:=\{v\in\R^d:\mu'v\ge0\text{ for every }\mu\in K\}.
\]
The following calculation explains the difference between the singleton-core
and multiple-core models.

\begin{lemma}[Dual-cone dichotomy]\label{lem:dual}
Suppose that $\cC$ and $\cR$ are nonempty.
\begin{enumerate}
\item If $\cC=\{1\}$, then
\begin{equation}\label{eq:dual-one-core}
\Lambda_1^*
=
\left\{v\in\R^d:
v_r\le0\ (r\in\cR),\quad
v_1+\sum_{r\in\cR}v_r\ge0
\right\}
=
\operatorname{cone}\{e_1,e_1-e_r:r\in\cR\}.
\end{equation}
\item If $|\cC|\ge2$, then
\begin{equation}\label{eq:dual-many-core}
\Lambda_1^*=\{0\}.
\end{equation}
\end{enumerate}
\end{lemma}

\begin{proof}
Suppose first that $\cC=\{1\}$.  Every $\mu\in\Lambda_1$ has
$\mu_1>0$, so conicity allows us to normalize $\mu_1=1$.  The inequality
\[
v_1+\sum_{r\in\cR}\mu_rv_r\ge0
\qquad\text{for every }(\mu_r)_{r\in\cR}\le1
\]
forces $v_r\le0$ by sending one $\mu_r$ to $-\infty$.  Given these sign
restrictions, the left side is minimized by setting every $\mu_r=1$, which
gives $v_1+\sum_rv_r\ge0$.  The two conditions are also sufficient.  The
generator representation follows from
\[
v
=
\left(v_1+\sum_rv_r\right)e_1
+\sum_r(-v_r)(e_1-e_r).
\]

Now suppose that $|\cC|\ge2$ and $v\in\Lambda_1^*$.  Fix a core index
$j$ and choose another core index $c\ne j$.  For every $s\in\R$, the mean
$e_c+se_j$ belongs to $\Lambda_1$.  Hence
\[
v_c+sv_j\ge0\qquad\text{for every }s\in\R,
\]
which forces $v_j=0$.  Thus all core components of $v$ vanish.  For a
redundant index $r$, fix a core index $c$ and consider
\[
\mu=e_c+se_r,
\qquad s\le1.
\]
This mean belongs to $\Lambda_1$, so $sv_r\ge0$ for every $s\le1$.
Taking $s=1$ and $s=-1$ gives $v_r=0$.  This holds for every redundant
index, proving~\eqref{eq:dual-many-core}.
\end{proof}

Lemma~\ref{lem:dual} does not say that a dominator exists whenever
$\Lambda_1^*$ is nontrivial.  It identifies the first-order imbalances that
are compatible with dominance.  In the singleton-core model, calibrated
gated dominators exploit such imbalances for particular tests.  In the
multiple-core model, the dual cone is trivial and dominance instead forces a
first-moment identity.  This is the sense in which the dual cone organizes
the results below.

\section{Feasible support}

For a closed convex set $A$, let
\[
h_A(a):=\sup_{x\in A}a'x
\]
denote its support function.  If $A$ is nonempty and coordinatewise lower,
then $h_A(a)<\infty$ for a nonzero $a$ only if $a\ge0$.

Relative to $P_0=N_d(0,\Omega)$, the Gaussian likelihood ratio is
\begin{equation}\label{eq:lr}
\frac{dP_\mu}{dP_0}(x)
=
\exp\left\{\mu'\Omega^{-1}x
-\frac12\mu'\Omega^{-1}\mu\right\}.
\end{equation}
A sample-space normal $a$ generates the natural-parameter ray $ta$ and the
mean ray $t\Omega a$.  Define the feasible-normal set
\begin{equation}\label{eq:feasible}
\cF(\Omega,\Lambda)
:=
\{a\in\R^d:a\ge0,\ \Omega a\in\Lambda\}.
\end{equation}

\begin{lemma}\label{lem:feasible-alternative}
If $a\in\cF(\Omega,\Lambda)$ and $a\ne0$, then
$\Omega a\in\Lambda_1$.
\end{lemma}

\begin{proof}
Suppose instead that every core component of $\Omega a$ is nonpositive.
Because $\Omega a\in\Lambda$, every redundant component is then also
nonpositive.  Since $a\ge0$, this would imply
$a'\Omega a\le0$, contradicting positive definiteness.
\end{proof}

\begin{definition}[Feasible support]\label{def:feasible-support}
For a nice region $A$, let
\[
F_A
:=
\{a\in\cF(\Omega,\Lambda):a\ne0,\ h_A(a)<\infty\}.
\]
We say that $A$ is \emph{feasibly supported} if
\begin{equation}\label{eq:feasible-support}
A
=
\bigcap_{a\in F_A}\{x:a'x\le h_A(a)\}.
\end{equation}
For a full-dimensional polyhedron, this is equivalent to feasibility of
every essential facet normal.
\end{definition}

\begin{proposition}[Feasible support suffices]\label{prop:feasible-support}
If a nice region $A$ is feasibly supported, then $\phi_A$ is
$\alpha$-admissible against all level-$\alpha$ tests, including randomized
competitors.
\end{proposition}

\begin{proof}
Suppose that a level-$\alpha$ test $\psi$ weakly dominates $\phi_A$ and
set $h:=\psi-\phi_A$.  Since $0\in\Lambda_0$,
$\E_0\psi\le\alpha=\E_0\phi_A$.  For any $c\in\cC$, the ray
$te_c$ belongs to $\Lambda_1$ for $t>0$.  Dominance on this ray and
continuity as $t\downarrow0$ give the reverse inequality.  Therefore
\begin{equation}\label{eq:zero-mass-support}
\E_0h=0.
\end{equation}
Moreover,
\begin{equation}\label{eq:signs}
\{h>0\}\subseteq A,
\qquad
\{h<0\}\subseteq A^c.
\end{equation}

Fix $a\in F_A$.  Lemma~\ref{lem:feasible-alternative} and conicity imply
$t\Omega a\in\Lambda_1$ for every $t>0$.  Dominance and
\eqref{eq:lr} give
\begin{equation}\label{eq:tilt-positive}
\E_0[e^{ta'X}h(X)]\ge0
\qquad\text{for every }t>0.
\end{equation}
Write $h=h^+-h^-$.  Since $h^+$ is supported on
$A\subseteq\{a'X\le h_A(a)\}$, for every $\varepsilon>0$,
\begin{align*}
\E_0[e^{ta'X}h(X)]
&\le
e^{t h_A(a)}\E_0h^+\\
&\quad-
e^{t\{h_A(a)+\varepsilon\}}
\E_0\left[h^-\1\{a'X\ge h_A(a)+\varepsilon\}\right].
\end{align*}
If the last expectation were positive, the right side would be negative for
all sufficiently large $t$, contradicting~\eqref{eq:tilt-positive}.
Taking $\varepsilon=1/n$ shows that $h^-=0$ almost surely on
$\{a'X>h_A(a)\}$.

By~\eqref{eq:feasible-support}, the open set $A^c$ is the union of these
open halfspaces over $a\in F_A$.  A countable subfamily covers $A^c$.
Consequently $h^-=0$ almost surely on $A^c$, and~\eqref{eq:signs} gives
$h^-=0$ almost surely everywhere.  Equation~\eqref{eq:zero-mass-support}
then forces $h^+=0$ almost surely.  Gaussian shifts are mutually absolutely
continuous, so a strict power improvement at another mean is impossible.
\end{proof}

The condition in Proposition~\ref{prop:feasible-support} is weaker than
requiring every finite-support direction to be feasible.  The distinction
disappears when there is one core moment.

\begin{proposition}[Relation with the finite-normal condition]
\label{prop:relation}
Let
\[
\cB(A):=\{b\in\R^d:b\ne0,\ h_A(b)<\infty\}.
\]
Then $\cB(A)\subseteq\cF(\Omega,\Lambda)$ implies feasible support.  If
$|\cC|=1$, the converse also holds.
\end{proposition}

\begin{proof}
The first implication follows from the standard supporting-halfspace
representation of a closed convex set.  Suppose now that $|\cC|=1$.
Then $\Lambda$ is defined by linear inequalities, so
$\cF(\Omega,\Lambda)$ is a closed convex cone.  Assume that $A$ is
feasibly supported and let $b\in\cB(A)$.  If
$b\notin\cF(\Omega,\Lambda)$, separation from this closed convex cone
provides $v\in\R^d$ such that
\[
a'v\le0\quad\text{for all }a\in\cF(\Omega,\Lambda),
\qquad b'v>0.
\]
For any $x\in A$ and $t\ge0$, every halfspace in
\eqref{eq:feasible-support} contains $x+tv$.  Feasible support therefore
implies $x+tv\in A$.  But then $h_A(b)=\infty$, a contradiction.
Thus $b$ is feasible.
\end{proof}

When $|\cC|\ge2$, the first implication can be strict.  For example, take
$d=3$, $\Omega=I_3$, $\cC=\{1,2\}$, and $\cR=\{3\}$.  The region
\begin{equation}\label{eq:symmetric-gate-example}
A_g
=
\{x_1\le\kappa,\ x_2\le\kappa,\ x_1+x_3\le b,\ x_2+x_3\le b\}
\end{equation}
has feasible facet normals $e_1$, $e_2$, $e_1+e_3$, and $e_2+e_3$.
It is therefore feasibly supported.  Yet $(1,1,2)$ has finite support on
$A_g$ and is infeasible because $2>\max\{1,1\}$.  Proposition
\ref{prop:feasible-support} covers this region after $(\kappa,b)$ are
calibrated to exact size, whereas the finite-normal condition does not.

\section{First-moment rigidity with multiple cores}

The next lemma imposes no shape restriction on the test being dominated.

\begin{lemma}[First-order restriction on a dominator]\label{lem:first-order}
Let $\phi$ be a level-$\alpha$ test satisfying $\E_0\phi=\alpha$, let
$\psi$ be another level-$\alpha$ test, and suppose that
$\E_\mu\psi\ge\E_\mu\phi$ for every $\mu\in\Lambda_1$.  Define
\[
h:=\psi-\phi,
\qquad
q:=\E_0[Xh(X)].
\]
Then
\begin{equation}\label{eq:first-order-conclusion}
\E_0h=0,
\qquad
\Omega^{-1}q\in\Lambda_1^*.
\end{equation}
In particular, if $|\cC|\ge2$, then $q=0$.
\end{lemma}

\begin{proof}
Level control gives $\E_0h\le0$.  Dominance along $te_c\in\Lambda_1$
and continuity at zero give $\E_0h\ge0$, proving the first identity.

Fix $\mu^*\in\Lambda_1$.  Since $\Lambda_1$ is a cone,
$t\mu^*\in\Lambda_1$ for every $t>0$.  Hence
\[
f(t):=\E_{t\mu^*}h(X)\ge0
\qquad(t>0),
\qquad f(0)=0.
\]
The Gaussian moment-generating function and $|h|\le1$ justify
differentiation from the right at zero.  Equation~\eqref{eq:lr} gives
\[
0\le f'(0+)
=\mu^{*\prime}\Omega^{-1}q.
\]
Thus $\Omega^{-1}q\in\Lambda_1^*$.  When $|\cC|\ge2$,
Lemma~\ref{lem:dual} gives $\Omega^{-1}q=0$, and nonsingularity of
$\Omega$ gives $q=0$.
\end{proof}

First-moment rigidity allows one supporting halfspace to be arbitrary.

\begin{proposition}[Feasible support plus one halfspace]
\label{prop:one-extra}
Suppose that $|\cC|\ge2$ and that a nice region can be written as
\begin{equation}\label{eq:one-extra-representation}
A=A_f\cap\{x:u'x\le\kappa\},
\qquad u\ne0,
\end{equation}
where
\[
A_f
=
\bigcap_{a\in G}\{x:a'x\le\kappa_a\}
\]
for some $G\subseteq\cF(\Omega,\Lambda)\setminus\{0\}$.  Then
$\phi_A$ is $\alpha$-admissible against all level-$\alpha$ tests,
including randomized competitors.
\end{proposition}

\begin{proof}
Suppose that $\psi$ weakly dominates $\phi_A$ and put
$h:=\psi-\phi_A$.  Lemma~\ref{lem:first-order} gives
\begin{equation}\label{eq:mass-moment-balance}
\E_0h=0,
\qquad
\E_0[Xh(X)]=0.
\end{equation}
For every $a\in G$, the exponential-tilt argument in the proof of
Proposition~\ref{prop:feasible-support}, with $\kappa_a$ in place of
$h_A(a)$, eliminates $h^-$ on $\{a'X>\kappa_a\}$.  A countable-subcover
argument gives
\[
h^-=0\quad\text{almost surely on }A_f^c.
\]
Since $h^-$ is supported on $A^c$, any remaining negative part is supported
on $A_f\cap\{u'X>\kappa\}$, whereas $h^+$ is supported on
$A\subseteq\{u'X\le\kappa\}$.

Let $M:=\E_0h^+=\E_0h^-$, where equality follows from
\eqref{eq:mass-moment-balance}.  If $M>0$, then
\[
\E_0[u'Xh^+]\le\kappa M,
\qquad
\E_0[u'Xh^-]>\kappa M.
\]
The second inequality is strict because $u'X>\kappa$ on the support of
$h^-$ and $u'X$ has no atom at $\kappa$.  It follows that
$u'\E_0[Xh]<0$, contradicting~\eqref{eq:mass-moment-balance}.  Therefore
$M=0$, so $h=0$ almost surely under $P_0$.  Mutual absolute continuity
again rules out strict improvement at another mean.
\end{proof}

\section{Consequences when there are at least two cores}

The first corollary removes the feasible-normal requirement from every
one-sided test.

\begin{corollary}[Every nonnegative halfspace is admissible]
\label{cor:any-halfspace}
Suppose that $|\cC|\ge2$.  Let $a\ge0$, $a\ne0$, and define
\begin{equation}\label{eq:one-sided-test}
\phi_a(x)
:=
\1\left\{a'x>
z_{1-\alpha}\sqrt{a'\Omega a}\right\},
\qquad
z_{1-\alpha}:=\Phi^{-1}(1-\alpha).
\end{equation}
Then $\phi_a$ is $\alpha$-admissible.  The conclusion includes weight
vectors concentrated entirely on redundant coordinates.
\end{corollary}

\begin{proof}
The acceptance halfspace is nice.  Apply Proposition~\ref{prop:one-extra}
with $A_f=\R^d$ and $u=a$.
\end{proof}

Thus feasibility is not necessary for admissibility when $|\cC|\ge2$.
For example, if $r\in\cR$, then $e_r$ is never feasible because
\[
(\Omega e_r)_r=1>max_{c\in\cC}\Omega_{cr},
\]
where strict inequality follows from positive definiteness and unit
diagonal.  Nevertheless, the test that rejects for a sufficiently large
$X_r$ is admissible by Corollary~\ref{cor:any-halfspace}.

The second corollary applies Proposition~\ref{prop:one-extra} to a familiar
multiple-testing region.

\begin{corollary}[Rectangles with one redundant moment]
\label{cor:rectangles}
Suppose that $|\cC|\ge2$, $|\cR|=1$, and $\Omega$ is any
positive-definite correlation matrix.  Let
\begin{equation}\label{eq:rectangle}
A_{\bm\kappa}
:=
\bigcap_{j=1}^d\{x_j\le\kappa_j\},
\qquad \kappa_j\in\R,
\end{equation}
and suppose that $\Pp_0(A_{\bm\kappa})=1-\alpha$.  Then the associated
test is $\alpha$-admissible.
\end{corollary}

\begin{proof}
Write $\cR=\{r\}$.  Every core face has a feasible normal.  Indeed, for
$c\in\cC$,
\[
(\Omega e_c)_r=\Omega_{rc}<1=(\Omega e_c)_c
\le\max_{c'\in\cC}(\Omega e_c)_{c'}.
\]
The redundant face is the one arbitrary halfspace.  Apply
Proposition~\ref{prop:one-extra}.
\end{proof}

In particular, the equal-cutoff max test is admissible whenever there are at
least two core moments and exactly one redundant moment.  This conclusion is
valid in arbitrary dimension and for every positive-definite correlation
matrix.

\section{The singleton-core regime}

When $\cC=\{1\}$, the first-order restriction in
Lemma~\ref{lem:first-order} becomes
\begin{equation}\label{eq:one-core-q}
(\Omega^{-1}q)_r\le0,
\qquad
(\Omega^{-1}q)_1+
\sum_{r\in\cR}(\Omega^{-1}q)_r\ge0.
\end{equation}
These inequalities allow a nonzero first-moment signature.  They do not, by
themselves, construct a dominator or imply that a particular test is
inadmissible.

Feasible halfspaces remain admissible.  In particular, one can always build
an admissible halfspace that assigns positive weight to every redundant
moment.  Fix $b\in\R^d$ with $b_1=0$ and $b_r>0$ for every
$r\in\cR$, and set
\begin{equation}\label{eq:epsilon-normal}
a_\varepsilon=e_1+\varepsilon b.
\end{equation}
At $\varepsilon=0$,
\[
(\Omega e_1)_1=1,
\qquad
(\Omega e_1)_r=\Omega_{r1}<1.
\]
Continuity and finiteness of $\cR$ imply that
$\Omega a_\varepsilon\in\Lambda$ for all sufficiently small
$\varepsilon>0$.  Proposition~\ref{prop:feasible-support} therefore gives
the following result.

\begin{corollary}[An admissible augmented halfspace]
\label{cor:epsilon}
Suppose that $\cC=\{1\}$.  For every positive-definite correlation matrix
$\Omega$, there is an $\varepsilon>0$ such that the exact-size test based on
$a_\varepsilon'X$ is $\alpha$-admissible and places strictly positive
weight on every redundant coordinate.
\end{corollary}

The nontrivial cone in~\eqref{eq:dual-one-core} also permits strict
domination.  The next result gives a general identity-covariance example.

\begin{proposition}[Pairwise gated domination of a rectangle]
\label{prop:gated-rectangle}
Suppose that $\Omega=I_d$, $\cC=\{1\}$, and
$\cR=\{2,\ldots,d\}$.  Let $A_{\bm\kappa}$ be the finite rectangle in
\eqref{eq:rectangle}, calibrated so that
\[
\prod_{j=1}^d\Phi(\kappa_j)=1-\alpha.
\]
Fix $r\in\cR$, and choose the unique finite $b_r$ satisfying
\begin{equation}\label{eq:gate-calibration-general}
\Pp_0\{X_1\le\kappa_1,\ X_1+X_r\le b_r\}
=
\Phi(\kappa_1)\Phi(\kappa_r).
\end{equation}
Define
\begin{equation}\label{eq:gated-rectangle-general}
A_{\bm\kappa}^{(r)}
:=
\{x_1\le\kappa_1,\ x_1+x_r\le b_r\}
\cap
\bigcap_{j\notin\{1,r\}}\{x_j\le\kappa_j\}.
\end{equation}
Then the test with acceptance region $A_{\bm\kappa}^{(r)}$ has exact size
$\alpha$ and strictly dominates the rectangular test at every
$\mu\in\Lambda_1$.  Hence every finite rectangle is inadmissible in this
model.
\end{proposition}

\begin{proof}
Independence and~\eqref{eq:gate-calibration-general} give
$\Pp_0(A_{\bm\kappa}^{(r)})=\Pp_0(A_{\bm\kappa})=1-\alpha$.
The gated region is coordinatewise lower, so its test is level $\alpha$.

In the $(1,r)$ coordinates, the newly rejected and newly accepted sets are
\begin{align*}
G_r
&:=\{x_1\le\kappa_1,\ x_r\le\kappa_r,\ x_1+x_r>b_r\},\\
D_r
&:=\{x_1\le\kappa_1,\ x_r>\kappa_r,\ x_1+x_r\le b_r\}.
\end{align*}
Calibration gives $\Pp_0(G_r)=\Pp_0(D_r)>0$.  Every $g\in G_r$ and
$d\in D_r$ satisfy
\[
g_1+g_r>d_1+d_r,
\qquad
-g_r>-d_r.
\]
For $\mu\in\Lambda_1$, write
\[
u:=\mu_1>0,
\qquad
v:=\mu_1-\mu_r\ge0.
\]
The likelihood ratio in these two coordinates is proportional to
\[
\exp\{u(x_1+x_r)+v(-x_r)\}.
\]
It is therefore strictly larger at every point of $G_r$ than at every point
of $D_r$.  Equal positive null mass gives
$\Pp_\mu(G_r)>\Pp_\mu(D_r)$.  The remaining independent rectangular
constraints contribute the common positive factor
\[
\prod_{j\notin\{1,r\}}\Phi(\kappa_j-\mu_j),
\]
which proves strict domination.
\end{proof}

The contrast is sharp for the configurations covered by the two rectangle
results.  Under identity covariance, a finite rectangle is strictly dominated
when there is one core moment and any positive number of redundant moments.
With at least two core moments and one redundant moment, the same class of
tests is admissible for every positive-definite correlation matrix.  The
mechanism is the change from the nontrivial cone in
\eqref{eq:dual-one-core} to the first-moment identity implied by
\eqref{eq:dual-many-core}.

\section{Maintained submodels}

The preceding first-moment argument uses conicity of $\Lambda_1$.  If the
maintained alternative is replaced by a smaller set
$\widetilde\Lambda_1$, the same argument applies directly when that set is
a cone, or is star-shaped at the origin along all of its directions.  For an
arbitrary maintained submodel, the relevant object is instead its cone of
local directions at the origin.  For example, define
\[
T_0(\widetilde\Lambda_1)
:=
\left\{v:\text{there exist }t_n\downarrow0\text{ and }v_n\to v
\text{ with }t_nv_n\in\widetilde\Lambda_1\right\}.
\]
The expansion used in Lemma~\ref{lem:first-order} then gives
\[
\Omega^{-1}q\in T_0(\widetilde\Lambda_1)^*.
\]
Thus a global dual cone should not be used mechanically for a maintained set
that contains no local rays.

Restrictions such as the existence of a fixed core coordinate $c^*$ with
\[
\mu_{c^*}\ge0,
\qquad
\mu_r\le\mu_{c^*}
\]
over the maintained alternative ensure nonnegative likelihood-ratio
coefficients for a pairwise gate.  They can therefore support a global gated
domination argument even when the maintained set is not conic.  Such a
condition exhibits useful transfer directions and makes the corresponding
global dual nontrivial.  It is sufficient, but not necessary, for that
nontriviality.

\section{What remains open}

The results establish two general sufficient conditions for nice tests:
feasible support for every nonempty core set, and feasible support plus one
arbitrary halfspace when $|\cC|\ge2$.  They do not establish that every nice
test is admissible in the multiple-core model.

For a nice polyhedron with $|\cC|\ge2$, any inadmissible example must require
at least two essential infeasible facets.  Two natural candidates are:
\begin{enumerate}
\item a finite rectangle with at least two redundant coordinates; and
\item in the identity-covariance model with
$\cC=\{1,2\}$ and $\cR=\{3\}$, the family
\begin{equation}\label{eq:two-bad-family}
A(\lambda)
:=
\{x_1\le\kappa,\ x_2\le\kappa,
x_1+\lambda x_3\le b,
x_2+\lambda x_3\le b\},
\qquad \lambda>1,
\end{equation}
with $(\kappa,b)$ calibrated to exact size.
\end{enumerate}
The last family has two nonproportional infeasible facet normals.  At
$\lambda=1$ it reduces to the feasibly supported symmetric gate.  Neither
Proposition~\ref{prop:feasible-support} nor
Proposition~\ref{prop:one-extra} determines its admissibility for
$\lambda>1$.

\begin{conjecture}[Multiple-core admissibility]
If $|\cC|\ge2$, every nice test is $\alpha$-admissible.
\end{conjecture}

The conjecture is consistent with the trivial dual cone and with all results
proved here, but those facts do not prove it.  The two-infeasible-facet
families are the first useful hostile checks.

Finally, the feasible-normal condition is known to be necessary within the
maintained class only in the bivariate model with one core and one redundant
moment, including every bivariate correlation matrix.  No necessity claim is
made here for a singleton core with several redundant moments or for an
arbitrary maintained submodel.

\end{document}
